% Part: many-valued-logic
% Chapter: syntax-and-semantics
% Section: sublogics

\documentclass[../../../include/open-logic-section]{subfiles}

\begin{document}

\olfileid[pt]{mvl}{syn}{sub}

\olsection{Lógicas multivaloradas como sublógicas de~$\LogCL$}

As lógicas multivaloradas usuais são todas definidas por matrizes nas quais
o valor de uma função de verdade para argumentos em $\{\True, \False\}$
coincide com o das funções de verdade clássicas. Especificamente, nessas
lógicas, se $x \in \{\True, \False\}$, então $\tf{\lnot}[\Log L](x) =
\tf{\lnot}[\LogCL](x)$, e, para $\star$ igual a qualquer um de $\land$, $\lor$,
$\lif$, se $x, y \in \{\True, \False\}$, então $\tf{\star}[\Log L](x,y) =
\tf{\star}[\LogCL](x,y)$. Em outras palavras, as funções de verdade de
$\lnot$, $\land$, $\lor$, $\lif$, restritas a $\{\True,\False\}$, são
exatamente as funções de verdade clássicas.

\begin{prop}\ollabel{prop:mvl-cl}
  Suponha que uma lógica multivalorada~$\Log L$ contenha em sua linguagem os conectivos
  $\lnot$, $\land$, $\lor$, $\lif$, que $\True, \False \in
  V$, e que suas funções de verdade satisfaçam:
  \begin{enumerate}
    \item\ollabel{prop:not} $\tf{\lnot}[\Log L](x) = \tf{\lnot}[\LogCL](x)$ se $x =
    \True$ ou $x = \False$;
    \item\ollabel{prop:land} $\tf{\land}[\Log L](x,y) = \tf{\land}[\LogCL](x,y)$,
    \item\ollabel{prop:lor} $\tf{\lor}[\Log L](x,y) = \tf{\lor}[\LogCL](x,y)$,
    \item\ollabel{prop:lif} $\tf{\lif}[\Log L](x,y) = \tf{\lif}[\LogCL](x,y)$,
    se $x, y \in \{\True, \False\}$.
  \end{enumerate}
  Então, para qualquer valoração $\pAssign v$ em~$V$ tal que $\pAssign
  v(p) \in \{\True,\False\}$, $\pValue v[\Log L](!A) = \pValue
  v[\LogCL](!A)$.
\end{prop}

\begin{proof}
  Por indução em~$!A$.
  \begin{enumerate}
  \item Se $!A \ident p$ é atômica, temos $\pValue v[\Log L](!A) =
  \pAssign v(p) = \pValue
  v[\LogCL](!A)$.
  \item Se $!A \ident \lnot B$, temos
  \begin{align*}
    \pValue v[\Log L](!A) & = \tf{\lnot}[\Log L](\pValue v[\Log L](!B)) & &\text{por \olref[val]{defn:pValue}}\\
    & = \tf{\lnot}[\Log L](\pValue v[\LogCL](!B)) && \text{pela hipótese de indução}\\
    & = \tf{\lnot}[\LogCL](\pValue v[\LogCL](!B)) 
&&\text{pela hipótese \olref{prop:not},}\\
&&&\text{pois $\pValue v[\LogCL](!B) \in \{\True, \False\}$,}\\
& = \pValue v[\LogCL](!A) &&\text{por \olref[val]{defn:pValue}}.
\end{align*}
\item Se $!A \ident (!B \land !C)$, temos
\begin{align*}
  \pValue v[\Log L](!A) & = \tf{\land}[\Log L](\pValue v[\Log L](!B), \pValue v[\Log L](!C)) & &\text{por \olref[val]{defn:pValue}}\\
  & = \tf{\land}[\Log L](\pValue v[\LogCL](!B),\pValue v[\LogCL](!C)) && \text{pela hipótese de indução}\\
  & = \tf{\land}[\LogCL](\pValue v[\LogCL](!B),\pValue v[\LogCL](!C)) 
&&\text{pela hipótese \olref{prop:land},}\\
&&&\text{pois $\pValue v[\LogCL](!B),\pValue v[\LogCL](!C) \in \{\True, \False\}$,}\\
& = \pValue v[\LogCL](!A) &&\text{por \olref[val]{defn:pValue}}.
\end{align*}
\end{enumerate}
Os casos em que $!A \ident (!B \lor !C)$ e $!A \ident (!B \lif !C)$
são semelhantes.
\end{proof}

\begin{cor}
  Se uma lógica multivalorada satisfaz as condições de
  \olref{prop:mvl-cl}, $\True \in V^+$ e $\False \notin V^+$, então
  ${\Entails[\Log L]} \subseteq {\Entails[\LogCL]}$, isto é, se $\Gamma
  \Entails[\Log L] !B$, então $\Gamma \Entails[\LogCL] !B$. Em
  particular, toda tautologia de $\Log L$ também é uma tautologia clássica.
\end{cor}

\begin{proof}
  Provamos a contrapositiva. Suponha $\Gamma \Entails/[\LogCL] !B$.
  Então existe alguma !!{valuation}~$\pAssign v\colon \PVar \to
  \{\True, \False\}$ tal que $\pValue v[\LogCL](!A) = \True$ para todo
  $!A \in \Gamma$ e $\pValue v[\LogCL](!B) = \False$. Como $\True,
  \False \in V$, a !!{valuation}~$\pAssign v$ também é !!a{valuation}
  para~$\Log L$. Por \olref{prop:mvl-cl}, $\pValue v[\Log L](!A) =
  \True$ para todo $!A \in \Gamma$ e $\pValue v[\Log L](!B) = \False$.
  Como $\True \in V^+$ e $\False \notin V^+$, isso significa que $\pAssign v
  \Entails[\Log L] \Gamma$ e $\pAssign v \Entails/[\Log L] !B$,
  isto é, $\Gamma \Entails/[\Log L] !B$.
\end{proof}
\end{document}
